\subsection{正则性的精确刻画}
\label{sec:ch14-precise-regularity-characterizations}
\setcounter{thmcount}{0}
\setcounter{equation}{0}

到目前为止, 第二个插值指标尚未得到重要使用. 不过, 它可以精确刻画多种正则性性质. 尖锐 Sobolev 不等式和迹定理就是例子(Scott 1977b). 对 $1<p<\infty$, $[L^p(\Omega),W_p^1(\Omega)]_{1/p,q}$ 中的函数在 $q=1$ 时都连续, 但在 $q>1$ 时并非如此; 相应的迹定理也成立(Scott 1977b).

另一个例子是, 分片光滑函数(Scott 1979)属于空间 $[L^2(\Omega),H^1(\Omega)]_{1/2,\infty}$, 但不属于任何更小的空间(例如, 跨越内部边界有间断的函数不属于 $H^{1/2}$); 参见习题~\ref{exr:ch14-08}.

\MathBookSourcePage{388}
若这样的函数 $f$ 是第~\ref{chap:ch05-n-dimensional-variational-problems} 章所讨论变分问题的数据, 则解 $u$ 将属于相应的插值空间. 例如, 假设正则性估计
\MBSourceItem{1}
\begin{equation}
\MathBookFormulaID{formula-ch14-piecewise-data-regularity-assumption}
\label{eq:ch14-piecewise-data-regularity-assumption}
\|u\|_{H^{k+2m}(\Omega)}\leq C\|f\|_{H^k(\Omega)}
\end{equation}
对某个 $k>1/2$ 成立(参见 Dauge 1988), 其中 $u$ 解
\[
\MathBookFormulaID{formula-ch14-piecewise-data-variational-problem}
a(u,v)=(f,v)\qquad\forall v\in V,
\]
而 $a(\cdot,\cdot)$ 在 $V\subset H^m(\Omega)$ 上连续且强制. 回顾(参见习题~\MBUnresolvedReference{exr:ch02-source-2-x-9}{2.x.9})还有
$\|u\|_{H^m(\Omega)}\leq C\|f\|_{H^m(\Omega)'}\leq C\|f\|_{H^{-m}(\Omega)}$. 对解算子 $f\mapsto u$ 作插值, 得
\MBSourceItem{2}
\begin{equation}
\MathBookFormulaID{formula-ch14-interpolated-piecewise-data-regularity}
\label{eq:ch14-interpolated-piecewise-data-regularity}
\|u\|_{[H^{2m}(\Omega),H^{2m+1}(\Omega)]_{1/2,\infty}}
\leq C\|f\|_{[H^0(\Omega),H^1(\Omega)]_{1/2,\infty}}.
\end{equation}
把这个正则性结果与定理~\ref{thm:ch14-intermediate-regularity-error} 的证明技术结合, 得到下述结果.

\MBSourceItem{3}
\begin{Theorem}
\label{thm:ch14-piecewise-smooth-data-error}
假设 $f$ 分片光滑(Scott 1979), ~\eqref{eq:ch14-piecewise-data-regularity-assumption} 对某个 $k>1/2$ 成立, 且~\eqref{eq:ch14-smooth-solution-error-estimate} 对 $k\geq2m+1$ 成立. 则
\[
\MathBookFormulaID{formula-ch14-piecewise-smooth-data-error}
\|u-u_h\|_{H^m(\Omega)}\leq Ch^{m+1/2}.
\]
\end{Theorem}
